\documentclass[12pt]{article}
\usepackage[utf8]{inputenc}
\usepackage[T1]{fontenc}
\usepackage{hyperref}
\title{Tangko Language Resources for K-12 Curriculum Development}
\author{Compiled by Claude Sonnet 5 (claude-sonnet-5) for Andr\'as Kornai}
\date{}
\begin{document}
\maketitle

\section{Introduction}
Tangko (ISO 639-3 tkx) is an Ok language of the Trans-New Guinea family, spoken in the village of Kawemaot in the Papua highlands of Indonesia. Sources give about 100 speakers (a 2000 figure). Sources disagree on the administrative district. One source says the language has never been written.

\section{Primary Educational Resources}
None found. No primers, textbooks or school programmes appeared in the search results.

\section{Linguistic and Reference Materials}
\begin{itemize}
\item Wikipedia overview: \url{https://en.wikipedia.org/wiki/Tangko_language}
\item Endangered Languages Project entry: \url{https://endangeredlanguages.com/lang/3539}
\item Numeral survey page (Wycliffe 2011 data; body-part counting up to 27): \url{https://lingweb.eva.mpg.de/channumerals/Tangko.htm}
\end{itemize}
No dictionary or grammar was found.

\section{Religious and Community Resources}
Nothing found.

\section{Digital and Audio Resources}
Nothing found.

\section{Curriculum Development Context}
Almost nothing is documented. Any curriculum work would need new documentation and an orthography developed with the community. The sources note that speakers in Kawemaot speak Indonesian poorly, which may matter for school planning.
\end{document}